\documentclass[12pt]{article}
\usepackage[utf8]{inputenc}
\usepackage[T1]{fontenc}
\usepackage[margin=1in]{geometry}
\usepackage{setspace}
\usepackage{xcolor}
\usepackage[colorlinks=true, urlcolor=blue!60!black, linkcolor=blue!60!black]{hyperref}
\setlength{\parindent}{0em}
\setlength{\parskip}{0.6em}
\renewcommand{\baselinestretch}{1.0}
\begin{document}
\begin{center}
{\LARGE \textit{Dred Scott v. Sandford}, 60 U.S. 393 (1857)\par}
\vspace{0.5em}
{\large Case Brief --- Two-Page Edition\par}
\end{center}

\section{Memory Jogger}
The Court held that people of African descent could not be U.S. citizens and struck down the Missouri Compromise, hastening the Civil War; the Thirteenth and Fourteenth Amendments repudiated it (\href{https://supreme.justia.com/cases/federal/us/60/393/}{Justia}).

\section{Detailed Case Facts}
Dred Scott's enslaver, Army surgeon John Emerson, took him to free Illinois and to Fort Snelling, where the Missouri Compromise had ``forever prohibited'' slavery. After returning to Missouri and Emerson's death, Scott sued for his family's freedom.

\section{Procedural History}
A St.\ Louis jury freed Scott in 1850, but the Missouri Supreme Court reversed. Scott sued John Sanford of New York in federal court on diversity; the jury found for Sanford, and Scott brought a writ of error.

\section{Judicial Votes}
7--2 against Scott, with all nine Justices writing. Taney for the Court (joined fully by Wayne); Nelson, Grier, Daniel, Campbell, and Catron concurring; McLean and Curtis dissenting.

\section{Holding}
Black Americans descended from enslaved people, free or not, cannot be ``citizens'' able to sue in diversity. The Missouri Compromise was void because excluding slave property from a territory violates Fifth Amendment due process. Scott's status was governed by Missouri law on his return.

\section{Analysis of Opinions}
\textbf{Taney:} Founding-era opinion excluded Black people from ``We the People''; the Territory Clause reached only lands held in 1787; and the Fifth Amendment protected slave property in the territories. \textbf{Concurrences:} Nelson rested on Missouri law alone; Catron voided the Compromise on treaty and state-equality grounds. \textbf{Curtis (dissenting):} free Black men voted in five States at ratification, so they were among ``the people of the United States,'' and the territorial ruling was extrajudicial. \textbf{McLean (dissenting):} slavery is ``a mere municipal regulation.''

\section{Examples: Future Applications}
\textbf{Same-Side:} A free Black Massachusetts merchant's diversity suit in 1858 must be dismissed. \textbf{Opposite-Side:} A modern State denying citizenship to U.S.-born children of a racial minority violates the Fourteenth Amendment's Citizenship Clause. \textbf{Fence-Sitter:} A slaveholder passing through New York with enslaved people raised the \textit{Lemmon} problem Nelson reserved.

\section{Critique}
\textit{Dred Scott} heads the anticanon as the ``greatest self-inflicted wound,'' built on distorted history that Curtis refuted. Originalists and their critics both claim it: Taney used original-intent rhetoric, while Scalia and Thomas cite it as the origin of substantive due process.

\section{Key Quotations}
\begin{quote}
``They had for more than a century before been regarded as beings of an inferior order, and altogether unfit to associate with the white race, either in social or political relations; and so far inferior, that they had no rights which the white man was bound to respect; and that the negro might justly and lawfully be reduced to slavery for his benefit.''
\end{quote}
\noindent --- Chief Justice Roger B. Taney, Opinion of the Court.

\end{document}
